While chronic sleep loss slowly erodes health, acute, total sleep deprivation causes the mind to rapidly disintegrate. Historical studies on extreme sleep deprivation reveal a predictable progression into psychosis \citep{waters2018deprivation}:

\begin{itemize}
    \item \textbf{After 24 hours:} Pronounced fatigue, irritability, concentration lapses, and a sense of "depersonalization" or mental fog. Cognitive impairment is comparable to being legally drunk \citep{waters2018deprivation}.
    \item \textbf{Around 48 hours:} Serious neurological symptoms appear, including visual distortions, simple hallucinations (e.g., seeing flashes), and disordered, illogical thinking \citep{waters2018deprivation}.
    \item \textbf{Around 72+ hours:} The brain can slide into full-blown, acute psychosis. This state is often indistinguishable from delirium or schizophrenia, characterized by complex hallucinations (seeing and hearing things that are not there), paranoid delusions, and incoherent speech \citep{waters2018deprivation}.
\end{itemize}